%%
%% fall-lecture09.tex
%% 
%% Made by Alex Nelson <pqnelson@gmail.com>
%% Login   <alex@lisp>
%% 
%% Started on  2026-10-17T09:16:53-0700
%% Last update 2026-10-17T09:16:53-0700
%% 

\lecture{}

\subsection{Bernoulli's equation}

\begin{node}
Last time, we saw for compressible fluids:
\begin{subequations}
\begin{equation}
h+\frac{1}{2}u_{i}u_{i}+V=\mbox{constant along stream lines}.
\end{equation}
For incompressible fluids, we derived (by multiplying through by
density $\rho$):
\begin{equation}
p+\frac{1}{2}\rho u_{i}u_{i}+\rho V = \mbox{constant along stream lines}.
\end{equation}
\end{subequations}
\end{node}

\begin{example}
Suppose we have a tank of fluid which has a ``spout'' at the
bottom. It looks something like:
\begin{center}
\includegraphics{img/img.14}
\end{center}
We assume:
\begin{enumerate}
\item The fluid is a steady, inviscid, incompressible flow (so we can
  use Bernoulli's equation);
\item There is no heat source $Q=0$ or heat flow $q_{i}=0$;
\item The only body force is gravity.
\end{enumerate}
We also assume the vertical distance from the top of the fluid to the
spout is $h$. This suggests a natural coordinate system where the $z$
axis is parallel to the vertical walls, and the spout is located at
$z=0$ (so the top of the fluid initially is at $z=h$).

We further assume the pressure at the top of the fluid $p_{\text{top}}$ is
$1~\mathrm{atm}$, and the pressure of the fluid at the spout $p_{\text{spout}}$ is
initially $1~\mathrm{atm}$:
\begin{equation}
p_{\text{top}}(t_{0})=p_{\text{spout}}(t_{0})=1~\mathrm{atm},
\end{equation}
where $t_{0}$ is the initial time.

Then Bernoulli's principle allows us to conclude that
\begin{equation}
\underbrace{p_{\text{top}}+\frac{1}{2}\overbrace{u_{i}u_{i}}^{=0}+\rho gh}_{\text{top}}=\underbrace{p_{\text{spout}}+\frac{1}{2}u_{2}^{2}+\rho g 0}_{\text{spout}}
\end{equation}
which can be rearranged to give us
\begin{equation}
gh=\frac{1}{2}u_{2}^{2},
\end{equation}
or more simply
\begin{equation}
u_{2}=\pm\sqrt{2gh}.
\end{equation}
This should look familiar, since it's the velocity of a point particle
experiencing constant acceleration after traveling a distance $h$.
\end{example}

\subsection{Velocity Gradient Tensor}

\begin{definition}
We define the \define{Velocity Gradient Tensor} $a_{ij}:=\partial_{j}u_{i}$
which has 9 independent components in three-dimensions (and $n^{2}$
independent components in $n$-dimensions).
\end{definition}

\begin{node}
Suppose we have two fluid parcels located at nearby points $P$ and $Q$
at time $t_{0}$, separated by the displacement vector $\delta\vec{x}$
initially, with velocities $\vec{u}_{P}$ at $P$ and $\vec{u}_{Q}$ at
$Q$:
\begin{center}
\includegraphics{img/img.15}
\end{center}
Then translating $\vec{u}_{P}$ to the base point $Q$, we see that
\begin{equation}
\vec{u}_{P}-\vec{u}_{Q}\approx\tens{a}\,\delta\vec{x},
\end{equation}
the difference
between the two velocity vectors is (to first order) the velocity
gradient tensor multiplied [as a matrix] by the displacement vector.
This is all working with $\|\delta\vec{x}\|$ ``sufficiently small''.
\end{node}

\begin{node}
We can always decompose a square matrix as the sum of a symmetric
matrix with an antisymmetric matrix. We have
\begin{subequations}
  \begin{align}
a_{ij} &= \frac{1}{2}(a_{ij}+a_{ji}) + \frac{1}{2}(a_{ij}-a_{ji})\\
&= \frac{1}{2}(\partial_{j}u_{i}+\partial_{i}u_{j}) + \frac{1}{2}(\partial_{j}u_{i}-\partial_{i}u_{j})\\
&=S_{ij}+\Omega_{ij}
  \end{align}
\end{subequations}
where $S_{ij}:=\frac{1}{2}(\partial_{j}u_{i}+\partial_{i}u_{j})$ is
the \define{Rate of Strain Tensor}, and $\Omega_{ij}:=\frac{1}{2}(\partial_{j}u_{i}-\partial_{i}u_{j})$
is the \define{Rate of Rotation Tensor}.
\end{node}

\begin{remark}
Observe $S_{ij}$ is the rate of deformation of a fluid elemeent. So
$S_{ij}\,\delta x_{j}$ gives a purely straining motion in the velocity.
\end{remark}

\begin{remark}
Eigenvector decomposition of $S_{ij}$ gives hte principal axes.
If we had a fluid parcel which is a sphere, then $S_{ij}$ will produce
an ellipsoid using the eigenvectors of $S_{ij}$ as the axes, and look
like $\lambda_{1}x_{1}^{2}+\lambda_{2}x_{2}^{2}+\lambda_{3}x_{3}^{2}=(r')^{2}$
where $x_{1}$, $x_{2}$, $x_{3}$ are taken to be the points in the
initial sphere.
\end{remark}

\begin{remark}
We see the trace of the rate-of-strain tensor is
\begin{subequations}
  \begin{align}
    \tr(S_{ij}) &= S_{11}+S_{22}+S_{33}\\
    &= S_{kk}\\
&= \partial_{k}u_{k} = \vec{\nabla}\cdot\vec{u}.
  \end{align}
\end{subequations}
For incompressible flows, this is zero, so
\begin{equation}
\tr(S_{ij})=0\quad\mbox{for incompressible flows}.
\end{equation}
\end{remark}

\begin{node}[Rotation matrix and vorticity]
The rotation matrix $\overline{\overline{\Omega}}$ has components
\begin{subequations}
  \begin{align}
\Omega_{ij}&=\frac{1}{2}\begin{pmatrix}
0 & (\partial_{2}u_{1}-\partial_{1}u_{2}) & (\partial_{3}u_{1}-\partial_{1}u_{3})\\
(\partial_{1}u_{2}-\partial_{2}u_{1}) & 0 & (\partial_{3}u_{2}-\partial_{2}u_{3})\\
(\partial_{1}u_{3}-\partial_{3}u_{1}) & (\partial_{2}u_{3}-\partial_{3}u_{2}) & 0
\end{pmatrix}\\
&=\frac{1}{2}\begin{pmatrix}
0 & -(\partial_{1}u_{2}-\partial_{2}u_{1}) & (\partial_{3}u_{1}-\partial_{1}u_{3})\\
(\partial_{1}u_{2}-\partial_{2}u_{1}) & 0 & -(\partial_{2}u_{3}-\partial_{3}u_{2})\\
-(\partial_{3}u_{1}-\partial_{1}u_{3}) & (\partial_{2}u_{3}-\partial_{3}u_{2}) & 0
\end{pmatrix}
  \end{align}
\end{subequations}
which has 3 independent components. We see these are just the
components of the vorticity
\begin{equation}
\vec{\omega}=\vec{\nabla}\times\vec{u}.
\end{equation}
More explicitly:
\begin{equation}
\Omega_{ij} = \frac{1}{2}\begin{pmatrix}0 & -\omega_{3} & \omega_{2} \\
  \omega_{3} & 0 & -\omega_{1}\\
  -\omega_{2} & \omega_{1} & 0
\end{pmatrix}
\end{equation}
or
\begin{equation}
\Omega_{ij}=\frac{-1}{2}\varepsilon_{ijk}\omega_{k}.
\end{equation}
What happens to fluid velocity $\delta\vec{u}$ of nearby fluid parcels
separated by displacement vector $\delta\vec{x}$? Wee see
\begin{equation}
\Omega_{ij}\,\delta
x_{j}=\frac{-1}{2}\varepsilon_{ijk}\omega_{k}\,\delta x_{j}=\frac{1}{2}(\vec{\omega}\times\delta\vec{x})_{i},
\end{equation}
which describes rigid body rotation.
\end{node}

\begin{example}[Couette flow]\label{ex:fall:lec09:couette-flow}
Consider the ``planar Couette flow'':
We have two walls parallel to each other (separated by a distance $h$)
and fluid filling the gap between the walls. The top wall moves at
constant velocity $\vec{u}_{0}$, but the bottom walls remains stationary.
This means there is only shear stress.
\begin{center}
\includegraphics{img/img.16}
\end{center}
We can determine the pressure gradient, which is typically sketched
out using the blue arrows in the figure.

A fluid parcel located initially at position $(x_{0},y_{0})$ will have
velocity $u(y)=u_{0}y/h$ and $v=0$. So the $y$ position remains fixed
for fluid parcels. We see that the velocity gradient tensor is
\begin{equation}
\overline{\overline{\tens{a}}} = \begin{pmatrix}0 & u_{0}/h\\
0 & 0
\end{pmatrix} = \frac{1}{2}\frac{u_{0}}{h}\begin{pmatrix}0 & 1\\1 & 0
\end{pmatrix} + \frac{1}{2}\frac{u_{0}}{h}\begin{pmatrix}0 & 1\\-1 & 0
\end{pmatrix}
\end{equation}.
Then a fluid parcel which is initiall a box will be transformed into a
rhombus. The $\overline{\overline{\Omega}}$ will merely rotate the
box, the $\overline{\overline{\tens{S}}}$ will shear the box.
\end{example}

\begin{remark}
There is \emph{a lot} more which could be said regarding Couette
flows. We could spend the rest of our lives studying it, and still
have more to discover.
\end{remark}

\subsection{Constitutive relationship}

\begin{puzzle}
How does the rate of strain tensor $S_{ij}$ relate to the stress
tensor $\sigma_{ij}=-p\delta_{i,j}+\tau_{ij}$? (The answer to this
question is precisely ``by using a constitutive relationship'', but
let us spell out what this means.)
\end{puzzle}

\begin{node}\index{Newtonian fluids}
Newtonian fluids have a linear relationship between $\tau_{ij}$ and
the velocity. In this situation, there is a fourth-order tensor
$A_{ijkl}$ such that
\begin{subequations}
  \begin{align}
\tau_{ij} &= A_{ijkl}\partial_{l}u_{k}\\
&=A_{ijkl}(S_{kl}+\Omega_{kl}.
  \end{align}
\end{subequations}
\end{node}

\begin{node}
Assuming the fluid is isotropic (which is a fairly benign assumption:
most fluids we encounter in this class \emph{will} be
isotropic\footnote{An example of a non-isotropic fluid would be some
kind of polymer whose behaviour changes depending on its
orientation.}). Then $A_{ijkl}$ is an isotropic fourth-order tensor,
which means we can write it as
\begin{equation}
A_{ijkl}=c_{1}\delta_{ij}\delta_{kl}+c_{2}\delta_{ik}\delta_{jl}+c_{3}\delta_{il}\delta_{jk}
\end{equation}
for some coefficients $c_{1}$, $c_{2}$, $c_{3}$.
\end{node}

\begin{node}
Now, since $\tau_{ij}=\tau_{ji}$ is symmetric, we can infer that
\begin{equation}
A_{ijkl}=A_{jikl}
\end{equation}
is symmetric in its first two indices. This forces $c_{2}=c_{3}$. We
call this coefficient the \define{Coefficient of Viscosity} and denote it:
\begin{equation}
\mu:=c_{2}=c_{3},
\end{equation}
and the remaining coefficient is called (creatively) the
\define{Second Coefficient of Viscosity} and denote it:
\begin{equation}
\lambda:=c_{1}.
\end{equation}
Then we can write
\begin{equation}
A_{ijkl}=\lambda\delta_{ij}\delta_{kl}+\mu(\delta_{ik}\delta{jl}+\delta_{il}\delta_{jk}).
\end{equation}
\end{node}

\begin{node}
We also can see that $A_{ijkl}$ is symmetric in its \emph{last two} indices:
\begin{equation}
A_{ijkl}=A_{ijlk}.
\end{equation}
Then
\begin{equation}
A_{ijkl}=\lambda\delta_{ij}\delta_{kl}+2\mu\delta_{ik}\delta_{jl}.
\end{equation}
Now, let us see what this means for the stress tensor $\tau_{ij}$:
\begin{subequations}
  \begin{align}
\tau_{ij} &= A_{iijkl}\partial_{\ell}u_{k}\quad\mbox{but only
  symmetric part of $\partial_{\ell}u_{k}$ survives}\\
&=A_{ijkl}S_{kl}\\
&=\lambda\delta_{ij}\tr(S_{kl})+2\mu S_{ij}.
  \end{align}
\end{subequations}
\end{node}

\begin{remark}
We can derive the full expression from physical arguments alone. Our
derivation used linear algebraic arguments.
\end{remark}

\begin{remark}
The tensor $\tau_{ij}$ has the same principal axes as $S_{ij}$.
\end{remark}

\begin{remark}
Since the trace of the rate-of-strain tensor is the divergence of the
velocity vector
\begin{equation}
S_{kk}=\vec{\nabla}\cdot\vec{u},
\end{equation}
this means that trace of $\tau_{ij}$ is:
\begin{subequations}
  \begin{align}
\tau_{ii} &= 3\lambda\tr(S_{kl})+2\mu S_{ii}\\
&=(3\lambda+2\mu)(\vec{\nabla}\cdot\vec{u}),
  \end{align}
\end{subequations}
so the average
\begin{equation}
\frac{1}{3}\tau_{ii}=\frac{1}{3}\tr(\tau_{ij})=\left(\lambda+\frac{2}{3}\mu\right)(\vec{\nabla}\cdot\vec{u}).
\end{equation}
This motivates the quantity $(\lambda+\frac{2}{3}\mu)$ which is
defined to be the \define{Bulk Viscosity} or \define{Volume Viscosity}
which we denote by
\begin{equation}
\mu_{V} := \left(\lambda+\frac{2}{3}\mu\right).
\end{equation}
Then we see (plugging these quantities back into the definition of
$\sigma_{ij}$) that the full Cauchy stress tensor is
\begin{subequations}
  \begin{align}
\sigma_{ij} &= -p\delta_{i,j}+\tau_{ij}\\
&=-p\delta_{i,j}+A_{ijkl}S_{kl}\\
&=-p\delta_{i,j}+\lambda(\vec{\nabla}\cdot\vec{u})\delta_{i,j}+2\mu S_{i,j}\\
&=-\left[p -\left(\lambda+\frac{2}{3}\mu\right)(\vec{\nabla}\cdot\vec{u})\right]\delta_{i,j}+2\mu\left(S_{ij}-\frac{1}{3}(\vec{\nabla}\cdot\vec{u})\delta_{i,j}\right)
  \end{align}
\end{subequations}
where, in the last line, the first term in square brackets is called
the \define{Mechanical Pressure}
\begin{equation}
p_{\text{mech}} := p -\left(\lambda+\frac{2}{3}\mu\right)(\vec{\nabla}\cdot\vec{u})
\end{equation}
and we refer to $p$ as the \define{Thermodynamic Pressure}.
\end{remark}

\begin{remark}
For a more thorough discussion of this constitutive relation business, see~\cite[\S3.3]{batchelor1973introduction}.
For more about ``mechanical pressure'' versus ``thermodynamic
pressure'', see~\cite[\S2-4.3]{white1991viscous}.
\end{remark}

